\begin{pastexam}[label={px:ntu-106-esoe-ds-07}]{臺大工科海洋所 106 學年度 第 7 題}
（設計）

\textbf{(10 \%)} Suppose you have a deque $D$ containing the numbers (1,2,3,4,5,6,7,8), in this
order. Suppose further that you have an initially empty queue $Q$. Give a pseudo-code
description of a function that uses only $D$ and $Q$ (and no other variables or objects) and
results in $D$ storing the elements (1,2,3,5,4,6,7,8), in this order.
\end{pastexam}

\begin{pastsoln}
\textbf{想法。} 只能用 $D$、$Q$，不能用變數暫存元素，所以元素只能從一個容器直接搬到另一個
容器（例如 \proc{Add-Last}$(D, \proc{Remove-First}(D))$）。 雙端佇列把前端拿下來接到後端，
等於把整列\textbf{轉一格}；目標只是把 $4$、$5$ 對調，做法是：先轉到 $4$ 在最前面，把 $4$ 暫放
進 $Q$，讓 $5$ 先接到後面，再把 $4$ 從 $Q$ 接上，最後轉回原來的起點。 雙端佇列的前端取出、
尾端加入（資結書第~6.2.6~節）寫成 \proc{Remove-First}、\proc{Add-Last}，佇列的加入、取出寫成
\proc{Enqueue}、\proc{Dequeue}。

\begin{pseudo}{只用 $D$ 和 $Q$ 對調 $4$ 和 $5$}
\Procedure{Swap-Four-Five}{$D, Q$}
    \State \proc{Add-Last}$(D, \proc{Remove-First}(D))$ \Comment{搬 $1$}
    \State \proc{Add-Last}$(D, \proc{Remove-First}(D))$ \Comment{搬 $2$}
    \State \proc{Add-Last}$(D, \proc{Remove-First}(D))$ \Comment{搬 $3$}
    \State \proc{Enqueue}$(Q, \proc{Remove-First}(D))$ \Comment{$4$ 暫放進 $Q$}
    \State \proc{Add-Last}$(D, \proc{Remove-First}(D))$ \Comment{搬 $5$}
    \State \proc{Add-Last}$(D, \proc{Dequeue}(Q))$ \Comment{$4$ 接在 $5$ 後面}
    \State \proc{Add-Last}$(D, \proc{Remove-First}(D))$ \Comment{搬 $6$}
    \State \proc{Add-Last}$(D, \proc{Remove-First}(D))$ \Comment{搬 $7$}
    \State \proc{Add-Last}$(D, \proc{Remove-First}(D))$ \Comment{搬 $8$}
\EndProcedure
\end{pseudo}
重複的列沒有寫成迴圈，是因為迴圈需要計數變數，題目不允許。

\textbf{追蹤}（$D$ 由前到後，$Q$ 由前到後）：
\begin{center}
\begin{tabular}{cll}
\toprule
做完第幾行 & $D$ & $Q$ \\
\midrule
開始 & $(1, 2, 3, 4, 5, 6, 7, 8)$ & $()$ \\
3 & $(4, 5, 6, 7, 8, 1, 2, 3)$ & $()$ \\
4 & $(5, 6, 7, 8, 1, 2, 3)$ & $(4)$ \\
5 & $(6, 7, 8, 1, 2, 3, 5)$ & $(4)$ \\
6 & $(6, 7, 8, 1, 2, 3, 5, 4)$ & $()$ \\
9 & $(1, 2, 3, 5, 4, 6, 7, 8)$ & $()$ \\
\bottomrule
\end{tabular}
\end{center}

\textbf{正確性。} 第 1--3 行讓 $D$ 轉了三格，$4$ 到了最前面。 第 4--6 行把 $4$、$5$ 從前端
移到後端，但先接 $5$ 再接 $4$，其餘元素的循環順序沒變。 到這裡 $D$ 從前端算起少轉了三格
（$6, 7, 8$ 還在前面），第 7--9 行再轉三格，$1$ 回到最前面，得到 $(1, 2, 3, 5, 4, 6, 7, 8)$。
結束時 $Q$ 是空的。 共 $9$ 次搬移，每次 $O(1)$。

另一種寫法是把 $1, 2, 3$ 依序放進 $Q$、$4$ 先轉到 $D$ 的後端、$5$ 放進 $Q$、再用
\proc{Remove-Last} 把 $4$ 放進 $Q$，最後把 $6, 7, 8$ 放進 $Q$ 再全部倒回 $D$。 一樣正確，只是
要 $17$ 次搬移。
\end{pastsoln}
